\VerdierPhysicalPageAnchor{0063}%
\VerdierSemanticAnchor{ch1-1-6-5-continuation}
from which it follows immediately:

\VerdierSemanticAnchor{ch1-1-6-5-deduction}
\VerdierSemanticAnchor{ch1-1-6-5-beta-s-coboundary}
\begin{equation*}
  \alpha_s(\sigma_1,\sigma_2)\alpha(\sigma_1,\sigma_2)^{-1}
  =\beta_s(\sigma_2)\beta_s(\sigma_1+\sigma_2)^{-1}\beta_s(\sigma_1),
\end{equation*}

\VerdierSemanticAnchor{ch1-1-6-5-eq3}
\VerdierSemanticAnchor{ch1-1-6-5-eq3-beta}
\VerdierSemanticAnchor{ch1-1-6-5-eq3-basis}
\begin{equation}
  \begin{aligned}
    \beta_s(\sigma_1)
      &=(-1)^{\left(
          \sum_{\substack{i<j\\s(i)>s(j)}}p_i p_j
        \right)},\\
    \beta_s(e_i)&=1,
      \qquad \forall i\in[1,n].
  \end{aligned}
  \tag{1.6.5.3}
\end{equation}

\VerdierSemanticAnchor{ch1-1-6-5-application}
Consequently, by (1.6.2), every map
$\beta\colon\mathbb Z^n\longrightarrow\mathbb Z/2\mathbb Z$ such that

\VerdierSemanticAnchor{ch1-1-6-5-eq4}
\begin{equation}
  \alpha_s(\sigma_1,\sigma_2)\alpha(\sigma_1,\sigma_2)^{-1}
  =\beta(\sigma_2)\beta(\sigma_1+\sigma_2)^{-1}\beta(\sigma_1)
  \tag{1.6.5.4}
\end{equation}

\VerdierSemanticAnchor{ch1-1-6-5-eq5-intro}
is of the form

\VerdierSemanticAnchor{ch1-1-6-5-eq5}
\begin{equation}
  \beta=\beta_s\eta,
  \tag{1.6.5.5}
\end{equation}

\VerdierSemanticAnchor{ch1-1-6-5-closure}
where $\eta\colon\mathbb Z^n\longrightarrow\mathbb Z/2\mathbb Z$ is
the unique group homomorphism whose values on the
$e_j\in\mathbb Z^n$ are $\beta(e_j)$.

\VerdierSemanticAnchor{ch1-1-6-6}
\VerdierSemanticAnchor{ch1-1-6-6-product-setup}
\noindent\textbf{1.6.6.}\quad
Let $(C_i,\mathbb Z,T_i,c_i)$, $1\leq i\leq n$, be
$\mathbb Z$-categories, and let $(D,\mathbb Z,T,c)$ be an additive
$\mathbb Z$-category. Denote by $(C,\mathbb Z^n,T,c)$ the product
category of the $(C_i,\mathbb Z,T_i,c_i)$. The group $\mathbb Z^n$ is
naturally equipped with an ordered basis $e_1,\ldots,e_n$. It acts up
to isomorphism on $C$ by the formula (we use here the right-hand
notation introduced in (1.3.12)):

\VerdierSemanticAnchor{ch1-1-6-6-eq1}
\begin{equation}
  (X_1,\ldots,X_n)[\sigma]
  =(X_1[l_1],\ldots,X_n[l_n]),
  \qquad
  \sigma=\sum_{i=1}^{n}l_i e_i.
  \tag{1.6.6.1}
\end{equation}

\VerdierSemanticAnchor{ch1-1-6-6-phi-action}
Let $\varphi\colon\mathbb Z^n\longrightarrow\mathbb Z$ be the
homomorphism
$\varphi(\sum_{i=1}^{n}l_i e_i)=\sum_{i=1}^{n}l_i$.
Let $\mathbb Z^n$ act on $D$ through $\varphi$ (1.4.1). In this
situation, we shall always use the following conventions:

\VerdierSemanticAnchor{ch1-1-6-6-sign-rule-1}
\noindent\textbf{Sign rule I:}
\emph{Functors from $C$ to $D$ will always, unless expressly stated
otherwise, be $\mathbb Z^n$-functors twisted by the Koszul cocycle
relative to the ordered basis $e_1,\ldots,e_n$.}

\VerdierSemanticAnchor{ch1-1-6-6-sign-rule-2}
\noindent\textbf{Sign rule II:}
\emph{morphisms of functors will, unless expressly stated otherwise,
be morphisms of $\mathbb Z^n$-functors relative to the trivial
homomorphism
$\beta\colon\mathbb Z^n\longrightarrow\mathbb Z/2\mathbb Z$ (1.4.6).}
